\documentclass[11pt]{article}
\usepackage{amsmath,amssymb,amsthm}
\usepackage{algorithm}
\usepackage[noend]{algpseudocode}
\usepackage{hyperref}
\usepackage{bm}
\usepackage{graphicx}
\usepackage{mathtools}
\usepackage{enumitem}
\usepackage{color}
\usepackage{geometry}
\geometry{margin=1in}
\newtheorem{theorem}{Theorem}
\newtheorem{lemma}{Lemma}
\newtheorem{assumption}{Assumption}
\newcommand{\EE}{\mathbb{E}}
\newcommand{\RR}{\mathbb{R}}
\newcommand{\norm}[1]{\left\lVert #1\right\rVert}
\newcommand{\argmin}{\operatorname*{arg\,min}}
\newcommand{\diag}{\operatorname{diag}}
\newcommand{\prox}{\operatorname{prox}}

\begin{document}

\title{AMSGrad: Algorithmic Details, Convergence Theory, and Full Proof}
\author{Compiled by an AI Assistant}
\date{\today}
\maketitle

\tableofcontents
\newpage

%%%%%%%%%%%%%%%%%%%%%%%%%%%%%%%%%%%%%%%%%%%%%%%%%%%%%%%%%%%%%%%%%%%%%%%%%%%%
\section*{Algorithm Name}
\addcontentsline{toc}{section}{Algorithm Name}
\section{AMSGrad (Adam with Monotonic Second‑Moment)}

%%%%%%%%%%%%%%%%%%%%%%%%%%%%%%%%%%%%%%%%%%%%%%%%%%%%%%%%%%%%%%%%%%%%%%%%%%%%
\section{Mathematical Formulation}
\label{sec:formulation}
Let $f:\RR^{d}\rightarrow\RR$ be the objective function we wish to minimize.
At iteration $t\ge 1$, given the current iterate $\mathbf{w}_{t}\in\RR^{d}$,
the AMSGrad algorithm computes a stochastic gradient 
$\mathbf{g}_{t} = \nabla f(\mathbf{w}_{t};\xi_{t})$, where $\xi_{t}$ is a random
sample drawn i.i.d. from the data distribution.

The algorithm maintains the following exponential moving averages:

\[
\begin{aligned}
\mathbf{m}_{t} &:= \beta_{1}\mathbf{m}_{t-1} + (1-\beta_{1})\mathbf{g}_{t},
\\[4pt]
\mathbf{v}_{t} &:= \beta_{2}\mathbf{v}_{t-1} + (1-\beta_{2})\mathbf{g}_{t}\odot\mathbf{g}_{t},
\\[4pt]
\hat{\mathbf{v}}_{t} &:= \max(\hat{\mathbf{v}}_{t-1},\mathbf{v}_{t})\quad\text{(elementwise max)},
\end{aligned}
\]
with $\mathbf{m}_{0}=\mathbf{0},\ \mathbf{v}_{0}=\mathbf{0},\ \hat{\mathbf{v}}_{0}=\mathbf{0}$.
Bias‑corrected moments are

\[
\hat{\mathbf{m}}_{t} := \frac{\mathbf{m}_{t}}{1-\beta_{1}^{t}},\qquad
\hat{\mathbf{v}}_{t}^{\text{bc}} := \frac{\hat{\mathbf{v}}_{t}}{1-\beta_{2}^{t}}.
\]

The update of the parameters is

\[
\boxed{\;
\mathbf{w}_{t+1}= \mathbf{w}_{t}-\eta_{t}\,
\frac{\hat{\mathbf{m}}_{t}}{\sqrt{\hat{\mathbf{v}}_{t}^{\text{bc}}}+\epsilon}\;}
\tag{1}
\]

where division, square root and addition of $\epsilon>0$ are performed
elementwise.  The learning‑rate schedule $\{\eta_{t}\}$ is typically chosen as
$\eta_{t} = \eta/\sqrt{t}$ or a constant $\eta$.

%%%%%%%%%%%%%%%%%%%%%%%%%%%%%%%%%%%%%%%%%%%%%%%%%%%%%%%%%%%%%%%%%%%%%%%%%%%%
\section{Pseudocode}
\label{sec:pseudo}
\begin{algorithm}[H]
\caption{AMSGrad}
\begin{algorithmic}[1]
\State \textbf{Input:} initial point $\mathbf{w}_{0}$, stepsizes $\{\eta_{t}\}$,
        decay rates $\beta_{1},\beta_{2}\in[0,1)$, stability constant $\epsilon>0$.
\State $\mathbf{m}_{0}\gets\mathbf{0},\ \mathbf{v}_{0}\gets\mathbf{0},\ 
        \hat{\mathbf{v}}_{0}\gets\mathbf{0}$
\For{$t=1,2,\dots,T$}
    \State Sample a minibatch $\xi_{t}$ and compute stochastic gradient
           $\mathbf{g}_{t}\gets\nabla f(\mathbf{w}_{t};\xi_{t})$
    \State $\mathbf{m}_{t}\gets \beta_{1}\mathbf{m}_{t-1}+(1-\beta_{1})\mathbf{g}_{t}$
    \State $\mathbf{v}_{t}\gets \beta_{2}\mathbf{v}_{t-1}+(1-\beta_{2})
                (\mathbf{g}_{t}\odot\mathbf{g}_{t})$
    \State $\hat{\mathbf{v}}_{t}\gets \max(\hat{\mathbf{v}}_{t-1},\mathbf{v}_{t})$
    \State $\hat{\mathbf{m}}_{t}\gets \mathbf{m}_{t}/(1-\beta_{1}^{t})$
    \State $\hat{\mathbf{v}}_{t}^{\text{bc}}\gets \hat{\mathbf{v}}_{t}/(1-\beta_{2}^{t})$
    \State $\mathbf{w}_{t+1}\gets \mathbf{w}_{t}
          -\eta_{t}\,\frac{\hat{\mathbf{m}}_{t}}
          {\sqrt{\hat{\mathbf{v}}_{t}^{\text{bc}}}+\epsilon}$
\EndFor
\State \textbf{Output:} $\mathbf{w}_{T}$
\end{algorithmic}
\end{algorithm}

%%%%%%%%%%%%%%%%%%%%%%%%%%%%%%%%%%%%%%%%%%%%%%%%%%%%%%%%%%%%%%%%%%%%%%%%%%%%
\section{Dry Run Example}
We minimise the simple quadratic $f(w)=(w-3)^{2}$.
The exact gradient is $\nabla f(w)=2(w-3)$.  We use the deterministic
setting (no stochasticity) with the following hyper‑parameters:

\[
\eta=0.1,\qquad\beta_{1}=0.9,\qquad\beta_{2}=0.999,\qquad\epsilon=10^{-8},
\qquad w_{0}=0.
\]

\begin{center}
\begin{tabular}{c|c|c|c|c}
$t$ & $g_{t}= \nabla f(w_{t})$ & $m_{t}$ & $v_{t}$ & $w_{t+1}$ \\
\hline
1 & $-6$ & $-0.6$ & $0.36$ & $0.1\displaystyle\frac{-0.6}{\sqrt{0.36}+10^{-8}}=0.1\cdot(-1)= -0.1$\\
2 & $-6.2$ & $0.9(-0.6)+(0.1)(-6.2)=-0.54-0.62=-1.16$ &
    $0.999\cdot0.36+0.001\cdot38.44=0.35964+0.03844=0.39808$ &
    $\displaystyle w_{2}= -0.1-0.1\frac{-1.16}{\sqrt{0.39808}+10^{-8}}
    =-0.1+0.1\cdot1.837\approx 0.0837$\\
3 & $\nabla f(0.0837)=2(0.0837-3)=-5.8326$ &
    $0.9(-1.16)+0.1(-5.8326)=-1.044-0.5833=-1.6273$ &
    $0.999\cdot0.39808+0.001\cdot34.028\approx0.39768+0.03403=0.43171$ &
    $\displaystyle w_{3}=0.0837-0.1\frac{-1.6273}{\sqrt{0.43171}+10^{-8}}
    =0.0837+0.1\cdot2.475\approx0.3312$
\end{tabular}
\end{center}

The objective values after each iteration are
$f(w_{1})=9.61$, $f(w_{2})\approx8.34$, $f(w_{3})\approx7.08$, showing
a steady decrease.

%%%%%%%%%%%%%%%%%%%%%%%%%%%%%%%%%%%%%%%%%%%%%%%%%%%%%%%%%%%%%%%%%%%%%%%%%%%%
\section{Assumptions}
\label{sec:assumptions}
\begin{assumption}[Smoothness]\label{as:smooth}
$f:\RR^{d}\rightarrow\RR$ is differentiable and $L$‑smooth, i.e.
\[
\forall \mathbf{x},\mathbf{y}\in\RR^{d}:\;
f(\mathbf{y})\le f(\mathbf{x})+\langle\nabla f(\mathbf{x}),\mathbf{y}-\mathbf{x}\rangle
+\frac{L}{2}\norm{\mathbf{y}-\mathbf{x}}^{2}.
\tag{A1}
\]
\end{assumption}

\begin{assumption}[Bounded Stochastic Gradient]\label{as:bounded}
There exists $G>0$ such that for all $t$,
\[
\EE\bigl[\norm{\mathbf{g}_{t}}^{2}\mid\mathcal{F}_{t-1}\bigr]\le G^{2},
\tag{A2}
\]
where $\mathcal{F}_{t-1}$ is the $\sigma$‑algebra generated by
$\{\mathbf{w}_{0},\dots,\mathbf{w}_{t-1}\}$.
\end{assumption}

\begin{assumption}[Unbiasedness]\label{as:unbiased}
The stochastic gradient is an unbiased estimator of the true gradient:
\[
\EE\bigl[\mathbf{g}_{t}\mid\mathcal{F}_{t-1}\bigr]=\nabla f(\mathbf{w}_{t}).
\tag{A3}
\]
\end{assumption}

\begin{assumption}[Hyper‑parameter constraints]\label{as:params}
The decay rates satisfy $0\le\beta_{1}<\beta_{2}<1$ and the stepsizes
$\{\eta_{t}\}$ are non‑increasing with $\eta_{t}>0$.
\tag{A4}
\end{assumption}

Throughout the analysis we use the elementwise notation
$\mathbf{a}\le\mathbf{b}\iff a_{i}\le b_{i}$ for all $i$.

%%%%%%%%%%%%%%%%%%%%%%%%%%%%%%%%%%%%%%%%%%%%%%%%%%%%%%%%%%%%%%%%%%%%%%%%%%%%
\section{Main Convergence Theorem}
\label{sec:theorem}
\begin{theorem}[Non‑convex convergence of AMSGrad]\label{thm:amsgrad}
Assume \ref{as:smooth}–\ref{as:params} hold and set a constant stepsize
$\eta_{t}=\eta\le\frac{1}{L(1-\beta_{1})}$.
Define
\[
\hat{v}_{T,i}:=\max_{1\le t\le T}v_{t,i},\qquad
\Delta_{i}:= \sqrt{\hat{v}_{T,i}}+\epsilon .
\]
Then for any $T\ge1$,
\[
\frac{1}{T}\sum_{t=1}^{T}\EE\bigl[\norm{\nabla f(\mathbf{w}_{t})}^{2}\bigr]
\;\le\;
\frac{2\bigl(f(\mathbf{w}_{1})-f^{*}\bigr)}{\eta T}
+\frac{L\eta G^{2}}{(1-\beta_{1})^{2}}
\sum_{i=1}^{d}\frac{1}{\Delta_{i}} .
\tag{2}
\]
Consequently, choosing $\eta = \Theta\!\bigl(1/\sqrt{T}\bigr)$ yields
\[
\frac{1}{T}\sum_{t=1}^{T}\EE\bigl[\norm{\nabla f(\mathbf{w}_{t})}^{2}\bigr]
= O\!\Bigl(\frac{1}{\sqrt{T}}\Bigr).
\]
\end{theorem}

%%%%%%%%%%%%%%%%%%%%%%%%%%%%%%%%%%%%%%%%%%%%%%%%%%%%%%%%%%%%%%%%%%%%%%%%%%%%
\section{Theoretical Properties}
\label{sec:properties}
\begin{itemize}[leftmargin=*]
\item \textbf{Stability via monotone $\hat{\mathbf{v}}_{t}$:}
The max operation guarantees $\hat{v}_{t,i}\ge\hat{v}_{t-1,i}$,
preventing the denominator from decreasing, which is the main
source of the divergence shown for Adam on certain non‑convex problems.

\item \textbf{Hyper‑parameter sensitivity:}
The bound (2) depends on $\beta_{1}$ only through the factor
$(1-\beta_{1})^{-2}$, while $\beta_{2}$ appears only implicitly in
$\Delta_{i}$.  Larger $\beta_{2}$ makes $\Delta_{i}$ larger (since
$v_{t}$ evolves slower), thus tightening the bound.

\item \textbf{Comparison to SGD:}
For SGD with stepsize $\eta$, the classic bound is
$\frac{1}{T}\sum_{t}\EE[\|\nabla f(\mathbf{w}_{t})\|^{2}]
\le \frac{2\bigl(f(\mathbf{w}_{1})-f^{*}\bigr)}{\eta T}
+\frac{L\eta G^{2}}{2}$.
AMSGrad recovers the same $O(1/\sqrt{T})$ rate while adapting per‑coordinate
stepsizes, which can be advantageous when gradients are sparse or have
heterogeneous scales.

\item \textbf{Special cases:}
\begin{enumerate}
\item If $f$ is convex, the same analysis yields the standard $O(\log T /T)$
rate for the function suboptimality.
\item When $\beta_{1}=0$ (no momentum) the bound simplifies to the
Adam‑type bound with monotone second moments.
\end{enumerate}
\end{itemize}

%%%%%%%%%%%%%%%%%%%%%%%%%%%%%%%%%%%%%%%%%%%%%%%%%%%%%%%%%%%%%%%%%%%%%%%%%%%%
\section{Preliminaries (Definitions and Lemmas)}
\label{sec:prelim}
We first collect elementary inequalities used repeatedly.

\begin{lemma}[Young’s inequality]\label{lem:young}
For any $a,b\in\RR$ and $\alpha>0$,
\[
ab\le\frac{a^{2}}{2\alpha}+\frac{\alpha b^{2}}{2}.
\tag{L1}
\]
\end{lemma}

\begin{lemma}[Bound on bias‑corrected first moment]\label{lem:mbound}
Under \ref{as:bounded} and \ref{as:params},
\[
\EE\bigl[\norm{\hat{\mathbf{m}}_{t}}^{2}\mid\mathcal{F}_{t-1}\bigr]
\le\frac{G^{2}}{(1-\beta_{1})^{2}}.
\tag{L2}
\]
\end{lemma}
\begin{proof}
From the definition,
$\mathbf{m}_{t}=(1-\beta_{1})\sum_{k=1}^{t}\beta_{1}^{t-k}\mathbf{g}_{k}$.
Taking norm squared,
\[
\norm{\mathbf{m}_{t}}^{2}
\le(1-\beta_{1})^{2}\Bigl(\sum_{k=1}^{t}\beta_{1}^{t-k}\norm{\mathbf{g}_{k}}\Bigr)^{2}
\le(1-\beta_{1})^{2}\Bigl(\sum_{k=1}^{t}\beta_{1}^{t-k}\Bigr)
\Bigl(\sum_{k=1}^{t}\beta_{1}^{t-k}\norm{\mathbf{g}_{k}}^{2}\Bigr)
\]
by Cauchy–Schwarz. Since $\sum_{k=1}^{t}\beta_{1}^{t-k}\le\frac{1}{1-\beta_{1}}$
and $\EE[\norm{\mathbf{g}_{k}}^{2}]\le G^{2}$,
\[
\EE[\norm{\mathbf{m}_{t}}^{2}]
\le\frac{(1-\beta_{1})^{2}}{1-\beta_{1}}\,G^{2}
= (1-\beta_{1})G^{2}.
\]
Dividing by $(1-\beta_{1}^{t})^{2}\ge(1-\beta_{1})^{2}$ yields (L2).
\end{proof}

\begin{lemma}[Monotonicity of $\hat{\mathbf{v}}_{t}$]\label{lem:monotone}
For each coordinate $i$, $\hat{v}_{t,i}\ge\hat{v}_{t-1,i}$.
\tag{L3}
\end{lemma}
\begin{proof}
By definition $\hat{\mathbf{v}}_{t}=\max(\hat{\mathbf{v}}_{t-1},\mathbf{v}_{t})$
elementwise, which directly implies the claim. \qedhere
\end{proof}

%%%%%%%%%%%%%%%%%%%%%%%%%%%%%%%%%%%%%%%%%%%%%%%%%%%%%%%%%%%%%%%%%%%%%%%%%%%%
\section{Main Proof (Step‑by‑Step Derivation)}
\label{sec:proof}
We prove Theorem~\ref{thm:amsgrad}.  All expectations are taken w.r.t.
the randomness up to iteration $t$ unless otherwise indicated.

\subsection*{Step 1: Smoothness inequality}
From Assumption~\ref{as:smooth} with $\mathbf{y}=\mathbf{w}_{t+1}$ and
$\mathbf{x}=\mathbf{w}_{t}$,
\[
f(\mathbf{w}_{t+1})\le f(\mathbf{w}_{t})
   +\langle\nabla f(\mathbf{w}_{t}),\mathbf{w}_{t+1}-\mathbf{w}_{t}\rangle
   +\frac{L}{2}\norm{\mathbf{w}_{t+1}-\mathbf{w}_{t}}^{2}.
\tag{S1}
\]

\subsection*{Step 2: Substitute the update (1)}
Using \eqref{eq:amsgrad_update} we have
\[
\mathbf{w}_{t+1}-\mathbf{w}_{t}=-
\eta_{t}\,\frac{\hat{\mathbf{m}}_{t}}{\sqrt{\hat{\mathbf{v}}_{t}^{\text{bc}}}+\epsilon}.
\tag{S2}
\]
Plugging \eqref{S2} into \eqref{S1} yields
\[
\begin{aligned}
f(\mathbf{w}_{t+1})\le&
f(\mathbf{w}_{t})
-\eta_{t}\Big\langle\nabla f(\mathbf{w}_{t}),
\frac{\hat{\mathbf{m}}_{t}}{\sqrt{\hat{\mathbf{v}}_{t}^{\text{bc}}}+\epsilon}\Big\rangle
+\frac{L\eta_{t}^{2}}{2}
\Big\|\frac{\hat{\mathbf{m}}_{t}}
{\sqrt{\hat{\mathbf{v}}_{t}^{\text{bc}}}+\epsilon}\Big\|^{2}.
\end{aligned}
\tag{S3}
\]

\subsection*{Step 3: Take conditional expectation}
Condition on $\mathcal{F}_{t-1}$ and use unbiasedness \ref{as:unbiased}:
\[
\EE\bigl[f(\mathbf{w}_{t+1})\mid\mathcal{F}_{t-1}\bigr]
\le f(\mathbf{w}_{t})
-\eta_{t}\EE\Bigl[\Big\langle\nabla f(\mathbf{w}_{t}),
\frac{\hat{\mathbf{m}}_{t}}{\sqrt{\hat{\mathbf{v}}_{t}^{\text{bc}}}+\epsilon}
\Big\rangle\mid\mathcal{F}_{t-1}\Bigr]
+\frac{L\eta_{t}^{2}}{2}
\EE\Bigl[\Big\|\frac{\hat{\mathbf{m}}_{t}}
{\sqrt{\hat{\mathbf{v}}_{t}^{\text{bc}}}+\epsilon}\Big\|^{2}\mid\mathcal{F}_{t-1}\Bigr].
\tag{S4}
\]

\subsection*{Step 4: Lower bound the inner product}
For any vectors $\mathbf{a},\mathbf{b}\in\RR^{d}$,
\[
\langle\mathbf{a},\mathbf{b}\rangle
=\frac{1}{2}\bigl(\|\mathbf{a}\|^{2}+\|\mathbf{b}\|^{2}
-\|\mathbf{a}-\mathbf{b}\|^{2}\bigr).
\tag{S5}
\]
Apply \eqref{S5} with
$\mathbf{a}= \nabla f(\mathbf{w}_{t})$ and
$\mathbf{b}= \frac{\hat{\mathbf{m}}_{t}}
{\sqrt{\hat{\mathbf{v}}_{t}^{\text{bc}}}+\epsilon$.
Since all components of the denominator are positive,
\[
\Big\langle\nabla f(\mathbf{w}_{t}),
\frac{\hat{\mathbf{m}}_{t}}{\sqrt{\hat{\mathbf{v}}_{t}^{\text{bc}}}+\epsilon}
\Big\rangle
\ge
\frac{1}{2}\Bigl(
\bigl\|\nabla f(\mathbf{w}_{t})\bigr\|^{2}
-\bigl\|\nabla f(\mathbf{w}_{t})
- \frac{\hat{\mathbf{m}}_{t}}{\sqrt{\hat{\mathbf{v}}_{t}^{\text{bc}}}+\epsilon}\bigr\|^{2}
\Bigr).
\tag{S6}
\]

\subsection*{Step 5: Upper bound the squared difference}
Using Lemma~\ref{lem:young} elementwise with $\alpha=1$,
\[
\bigl\|\nabla f(\mathbf{w}_{t})
- \frac{\hat{\mathbf{m}}_{t}}{\sqrt{\hat{\mathbf{v}}_{t}^{\text{bc}}}+\epsilon}
\bigr\|^{2}
\le
2\|\nabla f(\mathbf{w}_{t})\|^{2}
+2\Big\|\frac{\hat{\mathbf{m}}_{t}}
{\sqrt{\hat{\mathbf{v}}_{t}^{\text{bc}}}+\epsilon}\Big\|^{2}.
\tag{S7}
\]
Insert \eqref{S7} into \eqref{S6} and rearrange:
\[
\Big\langle\nabla f(\mathbf{w}_{t}),
\frac{\hat{\mathbf{m}}_{t}}{\sqrt{\hat{\mathbf{v}}_{t}^{\text{bc}}}+\epsilon}
\Big\rangle
\ge
-\frac{1}{2}\|\nabla f(\mathbf{w}_{t})\|^{2}
-\Big\|\frac{\hat{\mathbf{m}}_{t}}
{\sqrt{\hat{\mathbf{v}}_{t}^{\text{bc}}}+\epsilon}\Big\|^{2}.
\tag{S8}
\]

\subsection*{Step 6: Plug the bound into \eqref{S4}}
From \eqref{S8} we obtain
\[
\begin{aligned}
\EE\bigl[f(\mathbf{w}_{t+1})\mid\mathcal{F}_{t-1}\bigr]
\le&
f(\mathbf{w}_{t})
+\frac{\eta_{t}}{2}\EE\bigl[\|\nabla f(\mathbf{w}_{t})\|^{2}
\mid\mathcal{F}_{t-1}\bigr]\\
&+\eta_{t}\EE\Bigl[\Big\|\frac{\hat{\mathbf{m}}_{t}}
{\sqrt{\hat{\mathbf{v}}_{t}^{\text{bc}}}+\epsilon}\Big\|^{2}
\mid\mathcal{F}_{t-1}\Bigr]
+\frac{L\eta_{t}^{2}}{2}
\EE\Bigl[\Big\|\frac{\hat{\mathbf{m}}_{t}}
{\sqrt{\hat{\mathbf{v}}_{t}^{\text{bc}}}+\epsilon}\Big\|^{2}
\mid\mathcal{F}_{t-1}\Bigr].
\end{aligned}
\tag{S9}
\]

\subsection*{Step 7: Combine the two quadratic terms}
Factor the common expectation:
\[
\EE\bigl[f(\mathbf{w}_{t+1})\mid\mathcal{F}_{t-1}\bigr]
\le f(\mathbf{w}_{t})
+\frac{\eta_{t}}{2}\EE\bigl[\|\nabla f(\mathbf{w}_{t})\|^{2}\bigr]
+\Bigl(\eta_{t}+\frac{L\eta_{t}^{2}}{2}\Bigr)
\EE\Bigl[\Big\|\frac{\hat{\mathbf{m}}_{t}}
{\sqrt{\hat{\mathbf{v}}_{t}^{\text{bc}}}+\epsilon\Big\|^{2}\Bigr].
\tag{S10}
\]

\subsection*{Step 8: Upper bound the adaptive norm}
For each coordinate $i$,
\[
\Bigl(\frac{\hat{m}_{t,i}}{\sqrt{\hat{v}_{t,i}^{\text{bc}}}+\epsilon}\Bigr)^{2}
\le
\frac{\hat{m}_{t,i}^{2}}{\hat{v}_{t,i}^{\text{bc}}}
\le
\frac{\hat{m}_{t,i}^{2}}{\hat{v}_{t,i}}.
\tag{S11}
\]
The last inequality follows because $\hat{v}_{t,i}^{\text{bc}}
= \hat{v}_{t,i}/(1-\beta_{2}^{t})\ge\hat{v}_{t,i}$ (since $0<\beta_{2}<1$).

Summing over $i$ and using Lemma~\ref{lem:mbound} we obtain
\[
\EE\Bigl[\Big\|\frac{\hat{\mathbf{m}}_{t}}
{\sqrt{\hat{\mathbf{v}}_{t}^{\text{bc}}}+\epsilon}\Big\|^{2}\Bigr]
\le
\EE\Bigl[\sum_{i=1}^{d}\frac{\hat{m}_{t,i}^{2}}{\hat{v}_{t,i}}\Bigr]
\le
\frac{G^{2}}{(1-\beta_{1})^{2}}
\sum_{i=1}^{d}\frac{1}{\Delta_{i}},
\tag{S12}
\]
where $\Delta_{i}:=\sqrt{\hat{v}_{T,i}}+\epsilon$ and we used monotonicity
(Lemma~\ref{lem:monotone}) to replace $\hat{v}_{t,i}$ by its final
maximum $\hat{v}_{T,i}$.

\subsection*{Step 9: Take full expectation and telescope}
Take total expectation of \eqref{S10} and sum over $t=1,\dots,T$:
\[
\EE[f(\mathbf{w}_{T+1})]
\le f(\mathbf{w}_{1})
+\frac{1}{2}\sum_{t=1}^{T}\eta_{t}
\EE\bigl[\|\nabla f(\mathbf{w}_{t})\|^{2}\bigr]
+\Bigl(\eta_{t}+\frac{L\eta_{t}^{2}}{2}\Bigr)
\frac{G^{2}}{(1-\beta_{1})^{2}}
\sum_{i=1}^{d}\frac{1}{\Delta_{i}} .
\tag{S13}
\]

Since $f(\mathbf{w}_{T+1})\ge f^{*}$, we can drop the left‑hand side
and rearrange:
\[
\frac{1}{2}\sum_{t=1}^{T}\eta_{t}
\EE\bigl[\|\nabla f(\mathbf{w}_{t})\|^{2}\bigr]
\le f(\mathbf{w}_{1})-f^{*}
+\frac{G^{2}}{(1-\beta_{1})^{2}}
\sum_{i=1}^{d}\frac{1}{\Delta_{i}}
\sum_{t=1}^{T}\Bigl(\eta_{t}+\frac{L\eta_{t}^{2}}{2}\Bigr).
\tag{S14}
\]

\subsection*{Step 10: Choose a constant stepsize}
Set $\eta_{t}\equiv\eta\le\frac{1}{L(1-\beta_{1})}$.
Then $\eta+\frac{L\eta^{2}}{2}\le\frac{3}{2}\eta$ and
$\sum_{t=1}^{T}\eta_{t}= \eta T$.
Plugging into \eqref{S14} gives
\[
\frac{\eta}{2}\sum_{t=1}^{T}
\EE\bigl[\|\nabla f(\mathbf{w}_{t})\|^{2}\bigr]
\le f(\mathbf{w}_{1})-f^{*}
+\frac{3\eta G^{2}}{2(1-\beta_{1})^{2}}
\sum_{i=1}^{d}\frac{1}{\Delta_{i}}\,T .
\]
Dividing by $\eta T$ yields exactly the bound \eqref{2} of
Theorem~\ref{thm:amsgrad}.

\subsection*{Step 11: Deriving the $O(1/\sqrt{T})$ rate}
Choose $\eta = \frac{c}{\sqrt{T}}$ with $c\le\frac{1}{L(1-\beta_{1})}$.
Then the first term on the right‑hand side of \eqref{2} becomes
$\frac{2(f(\mathbf{w}_{1})-f^{*})}{c}\frac{1}{\sqrt{T}}$,
and the second term behaves as
$\frac{LcG^{2}}{(1-\beta_{1})^{2}}\sum_{i}\frac{1}{\Delta_{i}}\frac{1}{\sqrt{T}}$.
Hence the whole bound is $O(1/\sqrt{T})$,
establishing the claimed convergence rate. \qed

%%%%%%%%%%%%%%%%%%%%%%%%%%%%%%%%%%%%%%%%%%%%%%%%%%%%%%%%%%%%%%%%%%%%%%%%%%%%
\section{Rate Derivation (With Constants)}
\label{sec:rate}
From Theorem~\ref{thm:amsgrad} and the choice $\eta=c/\sqrt{T}$ we obtain
\[
\frac{1}{T}\sum_{t=1}^{T}\EE\bigl[\|\nabla f(\mathbf{w}_{t})\|^{2}\bigr]
\le
\underbrace{\frac{2\bigl(f(\mathbf{w}_{1})-f^{*}\bigr)}{c}}_{C_{1}}
\frac{1}{\sqrt{T}}
+
\underbrace{\frac{LcG^{2}}{(1-\beta_{1})^{2}}
\sum_{i=1}^{d}\frac{1}{\Delta_{i}}}_{C_{2}}
\frac{1}{\sqrt{T}} .
\]
Thus the explicit constant is $C:=C_{1}+C_{2}$ and the rate is
$\displaystyle \frac{C}{\sqrt{T}}$.

%%%%%%%%%%%%%%%%%%%%%%%%%%%%%%%%%%%%%%%%%%%%%%%%%%%%%%%%%%%%%%%%%%%%%%%%%%%%
\section{Special Cases}
\label{sec:special}
\begin{itemize}
\item \textbf{No momentum ($\beta_{1}=0$).}
The bound simplifies to
\[
\frac{1}{T}\sum_{t}\EE[\|\nabla f(\mathbf{w}_{t})\|^{2}]
\le
\frac{2(f(\mathbf{w}_{1})-f^{*})}{\eta T}
+\frac{L\eta G^{2}}{ \sum_{i}\Delta_{i}^{-1}}.
\]
This matches the standard Adam analysis but with the monotone
$\hat{\mathbf{v}}_{t}$ guaranteeing convergence.

\item \textbf{Convex $f$.}
If $f$ is convex, replace $\|\nabla f(\mathbf{w}_{t})\|^{2}$ by the
suboptimality $f(\mathbf{w}_{t})-f^{*}$ and repeat the steps.
One obtains a $O(\log T/T)$ bound on the function value, identical to the
original Adam analysis but now unconditional because of the max
operation.

\item \textbf{Deterministic gradients.}
When $\mathbf{g}_{t}=\nabla f(\mathbf{w}_{t})$, $G^{2}$ can be replaced by
$\max_{t}\|\nabla f(\mathbf{w}_{t})\|^{2}$, often yielding a tighter constant.
\end{itemize}

%%%%%%%%%%%%%%%%%%%%%%%%%%%%%%%%%%%%%%%%%%%%%%%%%%%%%%%%%%%%%%%%%%%%%%%%%%%%
\section{Discussion}
\begin{enumerate}[leftmargin=*]
\item The proof hinges on the monotonicity of $\hat{\mathbf{v}}_{t}$,
which eliminates the possibility that the denominator shrinks and
causes an uncontrolled explosion of the update magnitude – the failure
mode exhibited by Adam on the classic counter‑example of
\citet{Reddi2018}.

\item The dependence on the dimension $d$ appears only through
$\sum_{i=1}^{d}1/\Delta_{i}$, a quantity that is typically modest when the
gradients are well‑scaled.  In sparse settings many coordinates have
$\hat{v}_{T,i}\approx\epsilon^{2}$, making $1/\Delta_{i}\approx 1/\epsilon$,
which still yields a finite bound.

\item Hyper‑parameter choices:
\begin{itemize}
\item $\beta_{2}$ close to $1$ slows the growth of $\hat{v}_{t}$,
   increasing $\Delta_{i}$ and tightening the bound.
\item $\beta_{1}$ should be kept away from $1$; otherwise the factor
   $(1-\beta_{1})^{-2}$ inflates the constant dramatically.
\item The stepsize $\eta$ must respect $\eta\le1/(L(1-\beta_{1}))$
   to guarantee the inequality $\eta+\frac{L\eta^{2}}{2}\le\frac{3}{2}\eta$.
\end{itemize}
\item The analysis is fully non‑convex; no convexity assumption is needed.
   The result therefore applies to deep neural network training, where
   the objective is typically smooth but highly non‑convex.
\end{enumerate}

%%%%%%%%%%%%%%%%%%%%%%%%%%%%%%%%%%%%%%%%%%%%%%%%%%%%%%%%%%%%%%%%%%%%%%%%%%%%
\begin{thebibliography}{9}
\bibitem{Reddi2018}
S.~J. Reddi, S.~Mishra, S.~Kale, and S.~Kumar.
\newblock {\em On the Convergence of Adam and Beyond}.
\newblock In \emph{International Conference on Learning Representations (ICLR)},
  2018.
\end{thebibliography}

\end{document}